\chapter*{Preface}
The third part of \emph{Order and Chance} goes toward research: limit theorems with rates and large deviations, Brownian
motion and stochastic calculus, algebraic number theory and class field theory, curves and modular forms over finite fields,
sieves and the distribution of primes, and finally the meeting of the Riemann zeta function with random matrices. This book
is where those six tomes are put to work together. Its three hundred problems are frontier problems in two senses: they sit
on the borders between the tomes, and many of them stop a few steps away from questions that are still open.

The questions sound simple. How many choices of signs make $\pm1\pm2\pm\dots\pm n$ vanish? Why does a random polynomial of
degree one hundred have about three real roots? Why does a sum of random signs weighted by powers of the golden ratio refuse
to have a density? How often does a random permutation have a cycle structure that looks like the prime factorisation of a
random integer? The answers combine the local limit theorem with Poisson summation, Gaussian conditioning with complex
analysis, Fourier analysis with algebraic integers, and they are often more beautiful than the questions. Chapters are
organised by meeting points, never by tomes, and the statements never say which tool applies: finding it is half of each
problem.

The format is that of the first two problem books. A few \textbf{$\star$} problems open each section, \textbf{$\star\star$}
problems make a good afternoon, \textbf{$\star\star\star$} problems may take a weekend, and the problems marked
\emph{Journey} are long walks in many steps to a famous theorem. Hints, kept apart at the end of each chapter, give a nudge
when the search stalls. Every solution starts with \emph{How to find it}, which tells how one might think of the approach,
continues with a complete \emph{Write-up}, names the \emph{Tools used} with their tomes and chapters, and often ends with
\emph{Where it leads}, where research results quoted for their interest are labelled as such and never used in a proof.
Everything that is used is proved here or in a named tome.

Each chapter has a companion Python notebook. It recomputes every number in the chapter, from the $55\,252$ lucky tickets to
the constant in Kac's formula, runs the experiments that the problems suggest, and checks the statements that no proof in
this book can reach, such as the statistics of the zeros of the zeta function.

The pleasure of Part~III is that the frontier is close. A problem that starts with dice or coins can end, three steps later,
next to a theorem proved in the last ten years or a conjecture nobody can prove. We have tried to make those three steps
visible.
